\section{Method}
\label{sec:method}

All work was done in Python with PyTorch Geometric \parencite{PyG} for graphs, RDKit \parencite{RDKit} for chemistry and scikit-learn \parencite{sklearn} for the baseline and the metrics. The notebook runs end to end in about \SI{2.5}{\minute} on CPU and is deterministic at the seeds reported (\ref{sec:repro}).

\subsection{Dataset and splits}

Every loaded graph was converted to a record holding its SMILES string, attribute matrices and label, then checked for count, non-finite labels, unparsable molecules, empty graphs and class balance. Duplicate detection cannot use the raw strings, because one molecule has many valid SMILES spellings; structures were re-serialised to canonical SMILES, and the number of distinct labels per duplicated molecule then separates harmless repetition from contradiction. Splitting operates on unique molecules rather than records, otherwise one molecule can appear in both training and test; molecules whose duplicates carry contradictory labels were dropped, since the label is unknowable and keeping either copy would invent a ground truth.

The random split is stratified on the label at 70:15:15 with seed 42. The scaffold split computes scaffolds with chirality excluded, so stereoisomers share one, then fills training, validation and test greedily from the largest group downwards, never dividing a group. Two cases need an explicit rule. \textbf{Acyclic molecules} return the empty scaffold and are treated as one group, so they all land in a single subset: they have no ring system to generalise across, and scattering them would leak the easiest cases into test. \textbf{Equal-sized groups} are shuffled with the fixed seed before the size sort; breaking ties by file position instead produces a single-class test set on this dataset (\cref{sec:res-splits}), on which ROC-AUC is undefined, and shuffling changes only the arbitrary order among groups of equal size. Both splits were verified for coverage, disjointness, zero shared canonical SMILES and shared-scaffold count, then sealed until the final evaluation.

\subsection{The two models}

The baseline is Morgan fingerprints --- radius 2, 2048 bits, binary, unscaled --- classified by logistic regression (lbfgs, \num{5000} iterations) with balanced class weights from the training subset only; weighting was preferred to resampling, which would either discard roughly 700 positives or merely duplicate the minority. The inverse regularisation strength $C$ was selected from $\{0.01, 0.1, 1, 10, 100\}$ by validation ROC-AUC, which is threshold-free, and the threshold then selected separately on the same subset by Youden's index. One $C$ and one threshold per split strategy were frozen before any test molecule was scored.

The network applies \cref{eq:atom-encoder} with $d = 64$, then three \texttt{GINEConv} layers implementing \cref{eq:gine}, each wrapping a perceptron ($64 \to 128 \to 64$) with learnable $\varepsilon$ and followed by batch normalisation \parencite{BatchNorm}, ReLU and dropout $0.3$; then \cref{eq:pool} and a $64 \to 64 \to 1$ head --- \num{67588} parameters. Three layers give every atom a receptive field of three bonds, covering a functional group and its surroundings without over-smoothing molecules whose median size is 23 atoms, and the encoders train end to end, so no categorical code is ever consumed as a number. The loss is \texttt{BCEWithLogitsLoss} with $\text{pos\_weight} \approx 0.31$ from training counts only, matching the baseline's weighting; the optimiser is Adam \parencite{Adam} at learning rate \num{1e-3}; batches hold 64 molecules; early stopping has patience 25 on validation ROC-AUC with the best checkpoint restored, under a 150-epoch ceiling. The criterion is ROC-AUC rather than validation loss, which rises from over-confidence while the ranking keeps improving and would halt training near epoch 8. Three seeds (0, 1, 2) per split strategy control initialisation, dropout masks and batch order while the split seed stays 42, so seed variation and split-strategy effect are separable.

\subsection{Evaluation and error analysis}

Each test subset was scored once: the baseline at its frozen $C$ and threshold, each network run at the threshold maximising Youden's index on its own validation predictions. Metrics are aggregated over the three seeds as mean $\pm$ standard deviation against a majority-class reference. The error analysis uses the run with the \emph{median} test ROC-AUC on the scaffold split; its two most confident false positives and two most confident false negatives were taken, and for each the nearest training molecule by Tanimoto similarity retrieved with that neighbour's label and prediction, the median nearest-neighbour similarity over the test subset giving the scale.
